% assets/w05-course-crab.svg — 해류 삼각형: 선수각 · 침로각 · 크랩각 · 횡류각
%
%   bash _tools/tikz2svg.sh assets/src/w05-course-crab.tex assets/w05-course-crab.svg
%
% 요점: 대지속도는 대수속도와 조류의 벡터 합이다. 그 합이 뱃머리에서 틀어진 각이
%       크랩각이고, 그래서 선수각만 제어하면 배가 경로에서 옆으로 밀린다.
%
% 좌표는 계산한 값이다. 화면 규약: TikZ x = 동, TikZ y = 북.
% 방위각은 북에서 시계방향, TikZ 각 = 90 - 방위각.
%
%   psi = 35 deg (선수각),  u_r = 2.0, v_r = 0.35 m/s (대수),  U_c = 0.7 m/s, beta_c = 110 deg
%   조류 방위는 beta_c 로 적는다. 모델(build_w05_models 의 EOM)이 읽는 변수 이름이 beta_c 다.
%   Fossen 원문은 같은 것을 gamma_c 로 적지만, 이 볼트는 코드 이름을 따른다.
%     U_r      = sqrt(2.0^2 + 0.35^2) = 2.0304 m/s
%     beta_r   = atan2(0.35, 2.0)     = 9.926 deg      -> chi_r = 44.926 deg
%     U_r(NED) = (1.43733 북, 1.43394 동)
%     U_c(NED) = (-0.23941 북, 0.65778 동)
%     U (NED)  = ( 1.19792 북, 2.09172 동)  -> U = 2.4105 m/s, chi = 60.195 deg
%     beta     = chi - psi = 25.195 deg,  u = U cos beta = 2.1812, v = U sin beta = 1.0261
%     확인: atan2(v, u) = 25.195 deg = beta
%
%   그림 축척 S = 22 mm per (m/s)  ->  U_r 44.67 mm, U_c 15.40 mm, U 53.03 mm
%     U_r 끝 = (31.54, 31.63),  U 끝 = (46.02, 26.36)
%
%   아래 작은 그림 축척 S2 = 10 mm per (m/s), 원점 (12,-40)
%     U 끝 = (32.92, -28.02),  u 의 발 = (24.51, -22.13)   (수직성 내적 0 확인)
%
% 호의 반지름 — 북에서 재는 것은 안쪽(11/16/21), 뱃머리에서 재는 것은 바깥(28/34).
% 그래야 두 무리가 섞이지 않는다.
%
% 고친 것 (2026-10-17). 이전 그림은 psi, chi, beta 셋만 있었고 조류가 점선 화살표로
% 곁들여져 있을 뿐이어서, 크랩각이 "어디서 생기는가" 가 그림에 없었다.
% 대수속도 · 조류 · 대지속도의 벡터 삼각형을 그려 넣었다.
% 2026-09-18 에 고친 것(호를 동이 아니라 북에서 재기)은 그대로 지켰다.

\documentclass[border=6pt]{standalone}
% capstone-style.tex — 이 볼트의 TikZ 그림이 공유하는 스타일
%
% 쓰는 법:  % capstone-style.tex — 이 볼트의 TikZ 그림이 공유하는 스타일
%
% 쓰는 법:  % capstone-style.tex — 이 볼트의 TikZ 그림이 공유하는 스타일
%
% 쓰는 법:  \input{capstone-style}   (assets/src/ 안에서)
% 빌드:     bash _tools/tikz2svg.sh assets/src/w02-node-topic.tex assets/w02-node-topic.svg
%
% 왜 TikZ 인가
%   손으로 SVG 좌표를 쓰면 화살촉·선 굵기·글자 크기가 그림마다 미묘하게 달라진다.
%   그것이 "자료가 어설퍼 보인다"의 정체다. 스타일을 한 번만 정하고 상대좌표로
%   배치하면 좌표를 고쳐도 그림이 무너지지 않는다.
%   수식은 본문과 같은 TeX 글꼴로 나온다.
%
% 한글
%   ko.TeX(DVI 경로)를 쓴다. XeLaTeX 은 TikZ 그래픽을 PDF special 로 내보내는데
%   dvisvgm 이 그것을 받으려면 Ghostscript < 10.01 이 필요하다. 이 PC 에는
%   10.07 이 깔려 있어 그 경로가 막힌다. latex -> DVI -> dvisvgm 이 유일하게
%   동작하는 조합이다.

\usepackage[hangul]{kotex}
\usepackage{tikz}
\usepackage{amsmath}
\usepackage{xcolor}
\usetikzlibrary{arrows.meta, positioning, calc, fit, backgrounds, decorations.pathmorphing}

% ---- 색 --------------------------------------------------------------------
% 강조는 하나만. 나머지는 검정. 흑백 인쇄에서도 읽혀야 한다.
\definecolor{ink}{HTML}{1A1A1A}
\definecolor{soft}{HTML}{666666}
\definecolor{accent}{HTML}{7C3AED}
\definecolor{warn}{HTML}{B91C1C}

% 계층 색 — Simulink 모델의 블록 색과 같은 규칙을 쓴다
\definecolor{guid}{HTML}{1D4ED8}    % 유도
\definecolor{ctrl}{HTML}{EA580C}    % 제어
\definecolor{thr}{HTML}{D97706}     % 추진기
\definecolor{plant}{HTML}{16A34A}   % 운동모델
\definecolor{nav}{HTML}{7C3AED}     % 항법 · 신호처리
\definecolor{pathc}{HTML}{0D9488}   % 계획 경로 · 항적
\definecolor{errc}{HTML}{C2410C}    % 오차
\definecolor{flow}{HTML}{0369A1}    % 조류 · 바람

% ---- 선과 화살촉 -----------------------------------------------------------
\tikzset{
  sig/.style   = {ink, line width=0.5pt, -{Stealth[length=4pt, width=3pt]}},
  wire/.style  = {ink, line width=0.5pt},
  fb/.style    = {ink, line width=0.5pt, densely dashed,
                  -{Stealth[length=4pt, width=3pt]}},
  asig/.style  = {accent, line width=0.5pt, densely dashed,
                  -{Stealth[length=4pt, width=3pt]}},
  % 블록 — 흰 바탕, 직각, 검은 테두리. 색은 테두리에만 준다
  blk/.style   = {draw=ink, line width=0.55pt, fill=white,
                  minimum width=18mm, minimum height=8.5mm, inner sep=2pt, font=\small},
  gblk/.style  = {blk, draw=guid},
  cblk/.style  = {blk, draw=ctrl},
  tblk/.style  = {blk, draw=thr},
  pblk/.style  = {blk, draw=plant},
  nblk/.style  = {blk, draw=nav},
  % 합산점 · 분기점
  sum/.style   = {draw=ink, line width=0.55pt, fill=white, circle,
                  inner sep=0pt, minimum size=4.6mm, font=\scriptsize},
  dot/.style   = {ink, fill=ink, circle, inner sep=0pt, minimum size=1.5mm},
  % 글자
  sname/.style = {font=\small, inner sep=1.5pt},
  note/.style  = {font=\scriptsize, soft, inner sep=1.5pt},
  wnote/.style = {font=\scriptsize, warn, inner sep=1.5pt},
  anote/.style = {font=\scriptsize, accent, inner sep=1.5pt},
  axis/.style  = {ink, line width=0.5pt},
  tick/.style  = {font=\scriptsize, soft},
  % 묶음 상자 — 서브시스템 표시
  grp/.style   = {draw=soft, line width=0.4pt, densely dotted, rounded corners=1mm,
                  inner sep=3mm},
}
   (assets/src/ 안에서)
% 빌드:     bash _tools/tikz2svg.sh assets/src/w02-node-topic.tex assets/w02-node-topic.svg
%
% 왜 TikZ 인가
%   손으로 SVG 좌표를 쓰면 화살촉·선 굵기·글자 크기가 그림마다 미묘하게 달라진다.
%   그것이 "자료가 어설퍼 보인다"의 정체다. 스타일을 한 번만 정하고 상대좌표로
%   배치하면 좌표를 고쳐도 그림이 무너지지 않는다.
%   수식은 본문과 같은 TeX 글꼴로 나온다.
%
% 한글
%   ko.TeX(DVI 경로)를 쓴다. XeLaTeX 은 TikZ 그래픽을 PDF special 로 내보내는데
%   dvisvgm 이 그것을 받으려면 Ghostscript < 10.01 이 필요하다. 이 PC 에는
%   10.07 이 깔려 있어 그 경로가 막힌다. latex -> DVI -> dvisvgm 이 유일하게
%   동작하는 조합이다.

\usepackage[hangul]{kotex}
\usepackage{tikz}
\usepackage{amsmath}
\usepackage{xcolor}
\usetikzlibrary{arrows.meta, positioning, calc, fit, backgrounds, decorations.pathmorphing}

% ---- 색 --------------------------------------------------------------------
% 강조는 하나만. 나머지는 검정. 흑백 인쇄에서도 읽혀야 한다.
\definecolor{ink}{HTML}{1A1A1A}
\definecolor{soft}{HTML}{666666}
\definecolor{accent}{HTML}{7C3AED}
\definecolor{warn}{HTML}{B91C1C}

% 계층 색 — Simulink 모델의 블록 색과 같은 규칙을 쓴다
\definecolor{guid}{HTML}{1D4ED8}    % 유도
\definecolor{ctrl}{HTML}{EA580C}    % 제어
\definecolor{thr}{HTML}{D97706}     % 추진기
\definecolor{plant}{HTML}{16A34A}   % 운동모델
\definecolor{nav}{HTML}{7C3AED}     % 항법 · 신호처리
\definecolor{pathc}{HTML}{0D9488}   % 계획 경로 · 항적
\definecolor{errc}{HTML}{C2410C}    % 오차
\definecolor{flow}{HTML}{0369A1}    % 조류 · 바람

% ---- 선과 화살촉 -----------------------------------------------------------
\tikzset{
  sig/.style   = {ink, line width=0.5pt, -{Stealth[length=4pt, width=3pt]}},
  wire/.style  = {ink, line width=0.5pt},
  fb/.style    = {ink, line width=0.5pt, densely dashed,
                  -{Stealth[length=4pt, width=3pt]}},
  asig/.style  = {accent, line width=0.5pt, densely dashed,
                  -{Stealth[length=4pt, width=3pt]}},
  % 블록 — 흰 바탕, 직각, 검은 테두리. 색은 테두리에만 준다
  blk/.style   = {draw=ink, line width=0.55pt, fill=white,
                  minimum width=18mm, minimum height=8.5mm, inner sep=2pt, font=\small},
  gblk/.style  = {blk, draw=guid},
  cblk/.style  = {blk, draw=ctrl},
  tblk/.style  = {blk, draw=thr},
  pblk/.style  = {blk, draw=plant},
  nblk/.style  = {blk, draw=nav},
  % 합산점 · 분기점
  sum/.style   = {draw=ink, line width=0.55pt, fill=white, circle,
                  inner sep=0pt, minimum size=4.6mm, font=\scriptsize},
  dot/.style   = {ink, fill=ink, circle, inner sep=0pt, minimum size=1.5mm},
  % 글자
  sname/.style = {font=\small, inner sep=1.5pt},
  note/.style  = {font=\scriptsize, soft, inner sep=1.5pt},
  wnote/.style = {font=\scriptsize, warn, inner sep=1.5pt},
  anote/.style = {font=\scriptsize, accent, inner sep=1.5pt},
  axis/.style  = {ink, line width=0.5pt},
  tick/.style  = {font=\scriptsize, soft},
  % 묶음 상자 — 서브시스템 표시
  grp/.style   = {draw=soft, line width=0.4pt, densely dotted, rounded corners=1mm,
                  inner sep=3mm},
}
   (assets/src/ 안에서)
% 빌드:     bash _tools/tikz2svg.sh assets/src/w02-node-topic.tex assets/w02-node-topic.svg
%
% 왜 TikZ 인가
%   손으로 SVG 좌표를 쓰면 화살촉·선 굵기·글자 크기가 그림마다 미묘하게 달라진다.
%   그것이 "자료가 어설퍼 보인다"의 정체다. 스타일을 한 번만 정하고 상대좌표로
%   배치하면 좌표를 고쳐도 그림이 무너지지 않는다.
%   수식은 본문과 같은 TeX 글꼴로 나온다.
%
% 한글
%   ko.TeX(DVI 경로)를 쓴다. XeLaTeX 은 TikZ 그래픽을 PDF special 로 내보내는데
%   dvisvgm 이 그것을 받으려면 Ghostscript < 10.01 이 필요하다. 이 PC 에는
%   10.07 이 깔려 있어 그 경로가 막힌다. latex -> DVI -> dvisvgm 이 유일하게
%   동작하는 조합이다.

\usepackage[hangul]{kotex}
\usepackage{tikz}
\usepackage{amsmath}
\usepackage{xcolor}
\usetikzlibrary{arrows.meta, positioning, calc, fit, backgrounds, decorations.pathmorphing}

% ---- 색 --------------------------------------------------------------------
% 강조는 하나만. 나머지는 검정. 흑백 인쇄에서도 읽혀야 한다.
\definecolor{ink}{HTML}{1A1A1A}
\definecolor{soft}{HTML}{666666}
\definecolor{accent}{HTML}{7C3AED}
\definecolor{warn}{HTML}{B91C1C}

% 계층 색 — Simulink 모델의 블록 색과 같은 규칙을 쓴다
\definecolor{guid}{HTML}{1D4ED8}    % 유도
\definecolor{ctrl}{HTML}{EA580C}    % 제어
\definecolor{thr}{HTML}{D97706}     % 추진기
\definecolor{plant}{HTML}{16A34A}   % 운동모델
\definecolor{nav}{HTML}{7C3AED}     % 항법 · 신호처리
\definecolor{pathc}{HTML}{0D9488}   % 계획 경로 · 항적
\definecolor{errc}{HTML}{C2410C}    % 오차
\definecolor{flow}{HTML}{0369A1}    % 조류 · 바람

% ---- 선과 화살촉 -----------------------------------------------------------
\tikzset{
  sig/.style   = {ink, line width=0.5pt, -{Stealth[length=4pt, width=3pt]}},
  wire/.style  = {ink, line width=0.5pt},
  fb/.style    = {ink, line width=0.5pt, densely dashed,
                  -{Stealth[length=4pt, width=3pt]}},
  asig/.style  = {accent, line width=0.5pt, densely dashed,
                  -{Stealth[length=4pt, width=3pt]}},
  % 블록 — 흰 바탕, 직각, 검은 테두리. 색은 테두리에만 준다
  blk/.style   = {draw=ink, line width=0.55pt, fill=white,
                  minimum width=18mm, minimum height=8.5mm, inner sep=2pt, font=\small},
  gblk/.style  = {blk, draw=guid},
  cblk/.style  = {blk, draw=ctrl},
  tblk/.style  = {blk, draw=thr},
  pblk/.style  = {blk, draw=plant},
  nblk/.style  = {blk, draw=nav},
  % 합산점 · 분기점
  sum/.style   = {draw=ink, line width=0.55pt, fill=white, circle,
                  inner sep=0pt, minimum size=4.6mm, font=\scriptsize},
  dot/.style   = {ink, fill=ink, circle, inner sep=0pt, minimum size=1.5mm},
  % 글자
  sname/.style = {font=\small, inner sep=1.5pt},
  note/.style  = {font=\scriptsize, soft, inner sep=1.5pt},
  wnote/.style = {font=\scriptsize, warn, inner sep=1.5pt},
  anote/.style = {font=\scriptsize, accent, inner sep=1.5pt},
  axis/.style  = {ink, line width=0.5pt},
  tick/.style  = {font=\scriptsize, soft},
  % 묶음 상자 — 서브시스템 표시
  grp/.style   = {draw=soft, line width=0.4pt, densely dotted, rounded corners=1mm,
                  inner sep=3mm},
}


\begin{document}
\begin{tikzpicture}[x=1mm, y=1mm]

\coordinate (o)    at (0,0);
\coordinate (urt)  at (31.54,31.63);     % 대수속도 U_r 의 끝
\coordinate (ut)   at (46.02,26.36);     % 대지속도 U 의 끝

% ---- 북 · 동 축 ------------------------------------------------------------
\draw[axis, densely dashed, -{Stealth[length=4pt,width=3pt]}] (o) -- (0,52);
\node[font=\small, anchor=south] at (0,53) {$x$ (북)};
\draw[axis, -{Stealth[length=4pt,width=3pt]}] (o) -- (50,0);
\node[font=\small, anchor=west] at (51,0) {$y$ (동)};

% ---- 선체 축 ---------------------------------------------------------------
\draw[axis, densely dashed, -{Stealth[length=4pt,width=3pt]}] (o) -- (21.80,31.13);
\node[font=\small, anchor=south east] at (21.5,32.0) {뱃머리 $x_b$};
\draw[axis, densely dashed, -{Stealth[length=4pt,width=3pt]}] (o) -- (7.37,-5.16);
\node[font=\small, anchor=north west] at (7.9,-5.6) {$y_b$};

% ---- 속도 삼각형 -----------------------------------------------------------
% 대수속도 — 물에 대한 속도
\draw[pathc, line width=1.0pt, -{Stealth[length=5pt,width=4pt]}] (o) -- (urt);
\node[font=\small, pathc, anchor=south east] at (30.6,32.9) {$\boldsymbol{U}_r$};

% 조류 — 대수속도의 끝에 이어 붙인다
\draw[thr, line width=1.0pt, -{Stealth[length=5pt,width=4pt]}] (urt) -- (ut);
\node[font=\small, thr, anchor=south west] at (40.4,31.0) {$\boldsymbol{U}_c$};

% 대지속도 — 둘의 합
\draw[flow, line width=1.0pt, -{Stealth[length=5pt,width=4pt]}] (o) -- (ut);
\node[font=\small, flow, anchor=west] at (47.2,24.6) {$\boldsymbol{U}$};

% ---- 조류 방위 beta_c — 대수속도 끝에서 북으로부터 잰다 ---------------------
\draw[axis, densely dashed, line width=0.4pt] (urt) -- ++(0,9);
\draw[thr, line width=0.4pt] ($(urt)+(0,6)$) arc (90:-20:6);
\node[note, thr] at (38.0,38.4) {$\beta_c$};

% ---- 북에서 재는 세 각 -----------------------------------------------------
\draw[axis]  (o) ++(0,11) arc (90:55:11);
\node[note]  at (4.3,13.3) {$\psi$};
\draw[pathc] (o) ++(0,16) arc (90:45.074:16);
\node[note, pathc] at (7.6,17.9) {$\chi_r$};
\draw[flow]  (o) ++(0,21) arc (90:29.805:21);
\node[note, flow] at (11.5,21.6) {$\chi$};

% ---- 뱃머리에서 재는 두 각 -------------------------------------------------
\draw[accent] (o) ++(16.060,22.936) arc (55:45.074:28);
\node[note, accent] at (19.9,23.7) {$\beta_r$};
\draw[errc]   (o) ++(19.502,27.851) arc (55:29.805:34);
\node[note, errc] at (29.9,25.1) {$\beta$};

% ---- 배 -------------------------------------------------------------------
\begin{scope}[rotate=-35]
  \fill[ink] (-1.9,-2.7) -- (-1.9,1.6) -- (0,3.8) -- (1.9,1.6) -- (1.9,-2.7) -- cycle;
\end{scope}

% =========================================================================
% 오른쪽 — 읽는 법과 식
% =========================================================================
\node[font=\small, anchor=west] at (68,52) {해류 삼각형};

\node[note, anchor=west]         at (68,45.5) {$\psi$ \ 선수각 — 뱃머리가 향한 방향};
\node[note, flow, anchor=west]   at (68,40.3) {$\chi$ \ 침로각 — 배가 \textbf{땅에 대해} 가는 방향};
\node[note, pathc, anchor=west]  at (68,35.1) {$\chi_r$ \ 대수 침로각 — 배가 \textbf{물에 대해} 가는 방향};
\node[note, errc, anchor=west]   at (68,29.9) {$\beta$ \ 크랩각 — 뱃머리와 진행 방향의 차이};
\node[note, accent, anchor=west] at (68,24.7) {$\beta_r$ \ 횡류각 — 뱃머리와 물살 방향의 차이};
\node[note, thr, anchor=west]    at (68,19.5) {$\beta_c$ \ 조류가 흘러가는 방위 — 북에서 잰다};

\draw[soft, line width=0.3pt] (68,15.6) -- (130,15.6);

\node[font=\scriptsize, anchor=west] at (68,11.0)
     {$\boldsymbol{U} = \boldsymbol{U}_r + \boldsymbol{U}_c$ \quad 대지속도 = 대수속도 $+$ 조류};
\node[font=\scriptsize, anchor=west] at (68,5.0)
     {$U = \sqrt{u^2+v^2}$, \quad $U_r = \sqrt{u_r^2+v_r^2}$, \quad $U_c = \sqrt{u_c^2+v_c^2}$};
\node[font=\scriptsize, anchor=west] at (68,-1.0)
     {$\chi = \psi + \beta$, \quad $\beta = \operatorname{atan2}(v,\ u)$};
\node[font=\scriptsize, anchor=west] at (68,-7.0)
     {$\chi_r = \psi + \beta_r$, \quad $\beta_r = \arcsin\!\left(v_r/U_r\right)$};

\draw[soft, line width=0.3pt] (68,-11.5) -- (130,-11.5);

\node[note, anchor=west, text width=62mm] at (68,-17.5)
     {무엇이 무엇을 재는가 — GNSS 는 $\chi$ 와 $U$, 컴퍼스는 $\psi$, ADCP 는 조류를 잰다.
      $\beta$ 는 재는 것이 아니라 $\chi - \psi$ 로 얻는다};
\node[wnote, anchor=west, text width=62mm] at (68,-28.5)
     {조류가 없으면 $\boldsymbol{U} = \boldsymbol{U}_r$ 이고 $\beta = \beta_r$ 이 된다.
      이 그림의 두 각이 갈라지는 원인은 조류 하나다};
\node[wnote, anchor=west, text width=62mm] at (68,-39.5)
     {선수각 $\psi$ 만 제어하면 배는 경로에서 옆으로 밀린 채 나아간다.
      조류에서 경로 이탈이 남는 이유다};

% =========================================================================
% 아래 — 선체 축으로 본 같은 속도 (u, v 가 어디서 나오는가)
% =========================================================================
\node[font=\small, anchor=west] at (0,-14) {선체 축으로 본 대지속도 — $u$ 와 $v$};

\coordinate (p)  at (12,-40);
\coordinate (pu) at (24.51,-22.13);      % u 의 발
\coordinate (pt) at (32.92,-28.02);      % U 의 끝

\draw[axis, densely dashed, -{Stealth[length=4pt,width=3pt]}] (p) -- (25.19,-21.16);
\node[font=\scriptsize, anchor=south] at (25.19,-20.6) {$x_b$};
\draw[axis, densely dashed, -{Stealth[length=4pt,width=3pt]}] (p) -- (25.11,-49.18);
\node[font=\scriptsize, anchor=north west] at (25.5,-49.4) {$y_b$};

\draw[soft, densely dotted, line width=0.5pt] (p) -- (pu);
\draw[soft, densely dotted, line width=0.5pt] (pu) -- (pt);
\draw[flow, line width=1.0pt, -{Stealth[length=5pt,width=4pt]}] (p) -- (pt);

% 직각 표시 — u 의 발에서 x_b 방향과 y_b 방향으로 2.4mm 씩
\draw[soft, line width=0.4pt]
     ($(pu)+(1.38,1.97)$) -- ($(pu)+(1.38,1.97)+(1.97,-1.38)$) -- ($(pu)+(1.97,-1.38)$);

\node[note, soft] at (15.8,-29.35) {$u$};
\node[note, soft] at (30.4,-22.6)  {$v$};
\node[font=\scriptsize, flow, anchor=west] at (33.6,-28.6) {$\boldsymbol{U}$};
\draw[errc] (p) ++(8.030,11.468) arc (55:29.805:14);
\node[note, errc] at (24.63,-28.62) {$\beta$};

\begin{scope}[shift={(12,-40)}, rotate=-35]
  \fill[ink] (-1.6,-2.3) -- (-1.6,1.3) -- (0,3.2) -- (1.6,1.3) -- (1.6,-2.3) -- cycle;
\end{scope}

\node[note, anchor=west, text width=26mm] at (39,-31)
     {$u$ 전후, $v$ 좌우. 둘 다 \textbf{대지} 기준. 물에 대한 $u_r$, $v_r$ 도 같은 축으로 잰다};

\end{tikzpicture}
\end{document}
